8. Introducción a los Riesgos
Explicar qué protege el plan y cómo se conecta con el alcance, arquitectura, datos, metodologías, planificación, calidad e innovaciones. Presentar brevemente la función del T-16 y de los anexos. Aquí corresponde una introducción al capítulo, sin repetir todo el método ni listar los riesgos.
8.1 Plan de riesgos
Esta sección establece las reglas que después se aplicarán al proyecto.
- 8.1.1 Objetivo, alcance y enfoque de gestión. Definir los objetivos protegidos: alcance, plazo, calidad, seguridad, continuidad y viabilidad. Delimitar las fases consideradas y explicar cómo se aplicará ISO 31000. Distinguir riesgo, problema materializado, supuesto y restricción. Por ejemplo, el congelamiento estacional es una restricción conocida; una cancelación imprevista por marejada es un riesgo. Incorporar un diagrama del proceso de gestión.
- 8.1.2 Roles, responsabilidades y escalamiento. Definir quién mantiene el registro, quién responde por cada riesgo, quién ejecuta las acciones y quién autoriza su aceptación o el uso de reservas. Establecer cuándo se escala al Jefe de Proyecto o al Comité. Las responsabilidades deben coincidir con SD6 y considerar que el cliente no tiene contraparte informática.
- 8.1.3 Criterios de evaluación y aceptación de riesgos. Definir escalas de probabilidad e impacto, niveles de prioridad y condiciones para aceptar cada nivel. Propongo cinco niveles, con criterios explícitos para plazo, continuidad, datos, seguridad y adopción. Separar riesgo inicial de riesgo residual. Los valores y umbrales serán decisiones metodológicas de Synaptix, no cifras exigidas por las bases.
- 8.1.4 Registro y trazabilidad de los riesgos. Definir identificadores estables y campos del registro: causa, evento, impacto, categoría, fuente, etapa, evaluación, respuesta, responsable, disparador y estado. Explicar cómo se relacionan con T-16 y cómo se agregarán posteriormente los vínculos a EDT, actividades y pruebas, sin inventar códigos del SD7.
- 8.1.5 Ciclo de revisión, seguimiento y comunicación. Establecer la revisión en cada Comité de Proyecto y ante cambios relevantes, incidentes o proximidad de ventanas críticas. Explicar cómo se actualiza, comunica, escala y cierra un riesgo, y qué evidencia se requiere. La frecuencia operativa adicional debe coordinarse con SD6.
8.2 Identificación y Análisis de Riesgos
Esta sección aplica el método anterior a la solución concreta de Synaptix.
- 8.2.1 Fuentes de riesgo y estructura de desglose de riesgos. Presentar la RBS, que organiza los riesgos técnicos, organizacionales, de proyecto, seguridad y operación. Mostrar su relación con las bases, restricciones y decisiones de diseño. Incluir obsolescencia, dependencia de proveedor, escalabilidad, ciberseguridad y disponibilidad de contrapartes.
- 8.2.2 Identificación de riesgos de la solución, del desarrollo y de la implantación. Formular cada riesgo como causa → evento incierto → impacto, indicando componente y etapa afectados. Analizar los riesgos particulares del puerto: adopción por armadores, apoderados y contratistas; conocimiento del sistema antiguo; corrosión y flexión; marejadas; ausencia de contraparte técnica. Evaluarlos según las tecnologías efectivamente propuestas.
- 8.2.3 Análisis cualitativo y priorización. Aplicar las escalas, justificar las valoraciones y presentar la matriz de probabilidad e impacto. Identificar los riesgos prioritarios y explicar por qué concentran la exposición. Mostrar una síntesis en el cuerpo y conservar el registro completo en T-16. No atribuir eficacia comprobada a controles que todavía no se ejecutan.
- 8.2.4 Análisis cuantitativo de los riesgos prioritarios. Mostrar supuestos, fórmulas y consecuencias medibles: horas de indisponibilidad, duración de recuperación, registros que requieren conciliación o jornadas canceladas. Propongo analizar modos de falla para riesgos técnicos y escenarios de plazo para implantación. El efecto sobre la ruta crítica y la fecha final necesitará SD7. Una matriz de puntajes, por sí sola, no completa esta sección.
8.3 Plan de Acción a Riesgos
Esta sección convierte el análisis en respuestas ejecutables.
- 8.3.1 Estrategias de respuesta. Justificar si cada riesgo se evita, reduce, transfiere o acepta. Explicar qué causa o consecuencia modifica la respuesta. Distinguir controles ya incluidos en la arquitectura de acciones adicionales, para no contar dos veces el mismo esfuerzo.
- 8.3.2 Acciones, responsables, plazos y disparadores. Definir prevención, señal de activación, contingencia, responsables, plazo y evidencia de eficacia. Ahora se pueden usar hitos lógicos como «antes del corte de migración»; luego deberán vincularse a actividades y fechas reales del SD7.
- 8.3.3 Justificación costo-beneficio de las respuestas. Comparar esfuerzo, recursos y duración de las alternativas con las consecuencias que permiten evitar. Fundamentar por qué la respuesta elegida es proporcional. En el capítulo técnico se presenta esa justificación sin revelar precios; la valorización monetaria corresponde a la Oferta Económica.
- 8.3.4 Reservas de contingencia y de gestión y vinculación con el plan de trabajo. Explicar cómo se derivan, quién autoriza su uso y cómo se controlan. Vincular las reservas de plazo con marejadas, ventanas de intervención y migración. Su incorporación efectiva al cronograma se comprueba cuando exista SD7; no corresponde fijar ahora un porcentaje arbitrario.
Documentos de apoyo
Propongo acompañar el capítulo con:
- T-16: registro oficial de riesgos.
- A8.1: fichas complementarias de riesgos y respuestas.
- A8.2: memorias de análisis cuantitativo.
El siguiente trabajo puede ser redactar 8.1 completo, incluyendo escalas y responsabilidades propuestas. La estructura ya contempla el contenido exigido; los cálculos y su integración posterior todavía deben desarrollarse.